\begin{register}{H}{{mcliccfg}}{0x20800000}\label{regdesc:MCLICCFG_REG}
\regfield{(reserved)}{26}{6}{{0x00000}}%
\regfield{NMBITS}{2}{4}{{0x0}}%
\regfield{MNLBITS}{4}{0}{{0x0}}%
\reglabel{Reset}\regnewline%
\begin{regdesc}
\begin{reglist}
\label{fielddesc:MCLICCFG_MNLBITS}\item [MNLBITS] 配置 8 位的 \hyperref[regdesc:MCLICINTCTL_REG_i]{clicintctl[i]} 寄存器中有几个高位用来表示对应机器模式中断的级别。该配置为全局配置，对 HP 核心有效。\\
有效取值范围为 0-8。如果写入超过 8 的值，会被截断 (clamp) 为 8。 \\
(R/W)
\label{fielddesc:MCLICCFG_NMBITS}\item [NMBITS] 配置 \hyperref[fielddesc:MCLICINTATTR_i_MODE]{clicintattr[i].MODE} 中有几位用来表示中断模式。该配置为全局配置，对 HP 核心有效。对于仅有机器模式的系统，此域硬连线为 0。这表示 \hyperref[fielddesc:MCLICINTATTR_i_MODE]{clicintattr[i].MODE} 域硬连线为 0x3，因此所有中断仅在机器模式下。(RO) \\
\end{reglist}
\end{regdesc}
\end{register}


\begin{register}{H}{{clicinfo}}{0x20800004}\label{regdesc:CLICINFO_REG}
\regfield{(reserved)}{7}{25}{{0x00}}%
\regfield{CLICINTCTLBITS}{4}{21}{{0x3}}%
\regfield{(reserved)}{8}{13}{{0x00}}%
\regfield{NUM\_INTERRUPTS}{13}{0}{{0x0030}}%
\reglabel{Reset}\regnewline%
\begin{regdesc}
\begin{reglist}
\label{fielddesc:CLICINFO_NUM_INTERRUPTS}\item [NUM\_INTERRUPTS] 表示此 CLIC 实现支持的中断数量。硬连线为 48。(RO)
\label{fielddesc:CLICINFO_CLICINTCTLBITS}\item [CLICINTCTLBITS] 表示在 8 位 \hyperref[regdesc:MCLICINTCTL_REG_i]{clicintctl[i]} 寄存器中实际可编程的高位数。硬连线为 0x3。(RO)
\end{reglist}
\end{regdesc}
\end{register}


\begin{register}{H}{{clicintip[i]}}{0x20801000 + 4*i}\label{regdesc:MCLICINTIP_REG_i}
\regfield{(reserved)}{7}{1}{{0x00}}%
\regfield{{IP[i]}}{1}{0}{{0}}%
\reglabel{Reset}\regnewline%
\begin{regdesc}
\begin{reglist}
\label{fielddesc:MCLICINTIP_i}\item [{IP[i]}] 表示第 i 个 CLIC 机器模式中断的挂起状态。 \\
0：第 i 个中断未挂起 \\
1：第 i 个中断挂起 \\
\begin{itemize}
    \item 当使用 \hyperref[fielddesc:MCLICINTATTR_i_TRIG]{clicintattr[i].TRIG} 将第 i 个中断配置为电平类型时，此位为只读，若要清除此位，必须清除中断源。
    \item 当使用 \hyperref[fielddesc:MCLICINTATTR_i_TRIG]{clicintattr[i].TRIG} 将第 i 个中断配置为边缘触发时，此位为可读可写，若需清除此位，向此寄存器写入 0 即可。
\end{itemize}
(R/W)
\end{reglist}
\end{regdesc}
\end{register}

\begin{register}{H}{{clicintie[i]}}{0x20801001 + 4*i}\label{regdesc:MCLICINTIE_REG_i}
\regfield{(reserved)}{7}{1}{{0x00}}%
\regfield{{IE[i]}}{1}{0}{{0}}%
\reglabel{Reset}\regnewline%
\begin{regdesc}
\begin{reglist}
\label{fielddesc:MCLICINTIE_i}\item [{IE[i]}] 配置是否使能第 i 个 CLIC 机器模式中断。\\
0：禁用 \\
1：启用 \\
(R/W) \\
\end{reglist}
\end{regdesc}
\end{register}

\begin{register}{H}{{clicintattr[i]}}{0x20801002 + 4*i}\label{regdesc:MCLICINTATTR_REG_i}
\regfield{MODE}{2}{6}{{0x3}}%
\regfield{(reserved)}{3}{3}{{0x0}}%
\regfield{TRIG}{2}{1}{{0x0}}%
\regfield{SHV}{1}{0}{{0x0}}%
\reglabel{Reset}\regnewline%
\begin{regdesc}
\begin{reglist}
\label{fielddesc:MCLICINTATTR_i_SHV}\item [{SHV}] 配置第 i 个 CLIC 机器模式中断的硬件向量化。 \\
0：第 i 个中断未硬件向量化。即，CPU 在接受此中断时将跳转到 \hyperref[fielddesc:MTVECBASE]{mtvec} 中配置的地址。 \\
1：第 i 个中断硬件向量化。即，CPU 在接受此中断时，跳转的地址为相对于 \hyperref[fielddesc:MTVTBASE]{mtvt} 基地址偏移 (mtvt + 4*i) 处存储的字地址。 \\
(R/W)
\label{fielddesc:MCLICINTATTR_i_TRIG}\item [{TRIG}] 配置第 i 个 CLIC 机器模式中断的触发类型和极性。 \\
0x0：中断为正电平触发 \\
0x1：中断为正边缘触发 \\
0x2：中断为负电平触发 \\
0x3：中断为负边缘触发 \\
(R/W)
\label{fielddesc:MCLICINTATTR_i_MODE}\item [{MODE}] 配置第 i 个 CLIC 中断的特权模式。 \\
此值硬连线为 0x3，因为不支持用户模式中断。\\
(RO)
\end{reglist}
\end{regdesc}
\end{register}

\begin{register}{H}{{clicintctl[i]}}{0x20801003 + 4*i}\label{regdesc:MCLICINTCTL_REG_i}
\regfield{{CTL[i]}}{3}{5}{{0x0}}%
\regfield{(reserved)}{5}{0}{{0x1f}}%
\reglabel{Reset}\regnewline%
\begin{regdesc}
\begin{reglist}
\label{fielddesc:MCLICINTCTL_i}\item [{CTL[i]}] 配置第 i 个 CLIC 机器模式中断的级别和优先级。 \\
\begin{itemize}
    \item 由于在此 CLIC 实现中 \hyperref[fielddesc:CLICINFO_CLICINTCTLBITS]{CLICINTCTLBITS} 参数为 3，因此在 \hyperref[regdesc:MCLICINTCTL_REG_i]{clicintctl[i]} 中实现的位数相同，因此此寄存器的低 5 位硬连线为 1。 \\
    \item \hyperref[fielddesc:MCLICCFG_MNLBITS]{mcliccfg.MNLBITS} 的配置值决定了 \hyperref[regdesc:MCLICINTCTL_REG_i]{clicintctl[i]} 中表示中断级别的高位数，而剩余的低位表示优先级。
\end{itemize}
(R/W)
\end{reglist}
\end{regdesc}
\end{register}
